\documentclass[11pt]{article}
\usepackage{amsmath,amssymb,mathtools}
\usepackage{fontspec}
\usepackage{unicode-math}
\setmainfont{Latin Modern Roman}
\setmathfont{Latin Modern Math}
\usepackage{xurl}
\usepackage[unicode,colorlinks=true,linkcolor=blue,urlcolor=blue,pdftitle={Problem statement}]{hyperref}
\newcommand{\sref}[2]{\hyperlink{source:#1:#2}{[S#2]}}
\newcommand{\cataloguescope}[1]{\paragraph{Mathematical statement.}#1}
\begin{document}
\section{Research Problems in Function Theory — Problem 2.31}\label{statement}

\subsection{Definitions and mathematical statement}

\paragraph{Shared notation.}$\mathbb N=\{1,2,\ldots\}$, and $\mathbb Z,\mathbb Q,\mathbb R,\mathbb C$ denote the integers, rationals, reals and complex numbers. Any other conventions are specified locally.


Let $A,B\subset\mathbb C$ be countable dense sets; density means that every nonempty open disc meets the set. For a function $f$, $f^{-1}(B)=\{z:f(z)\in B\}$.

\paragraph{Statement recorded in the supplied source.}
For every such pair $A,B$, does there exist an entire function $f$ such that
\[f^{-1}(B)=A?\]
If the universal assertion is false, determine conditions on $A,B$ under which it holds. \sref{statement}{2}
Update 2.31 records Barth and Schneider's affirmative construction. This question is about the inverse-image identity; it does not additionally require $f(A)=B$. \sref{statement}{2}
\subsection{Short English statement}
Can an entire function arrange that membership of its value in one prescribed dense countable set occurs exactly on another prescribed dense countable set?
\subsection{Sources}
\begin{description}
\item[\hypertarget{source:statement:1}{S1}] ULAM.AI and the UnsolvedMath Contributors, \emph{UnsolvedMath: A Curated Collection of Open Mathematics Problems}, catalogue entry. CC BY 4.0. \url{https://huggingface.co/datasets/ulamai/UnsolvedMath}.
\item[\hypertarget{source:statement:2}{S2}] W. K. Hayman and E. F. Lingham, \emph{Research Problems in Function Theory} (2018), Chapter 2, Problem 2.31 and Update 2.31. \url{https://arxiv.org/abs/1809.07200}.
\end{description}
\label{end:statement}
\end{document}
